\chapter*{前言}
\addcontentsline{toc}{chapter}{前言}

多体量子动力学的困难并不只在于 Hilbert 空间随自由度指数增长，还在于研究者通常面对一个更具体的工作问题：初态已经给定，哈密顿量也已经写出，究竟怎样把它们转化成可以运行、可以验证、并能回到实验观测量的计算流程？本书围绕这一问题组织 Wigner 相空间方法与截断维格纳近似。

本书不把相空间方法描述成将量子力学“变成经典力学”的魔法。Wigner 函数是准概率分布，负值、算符排序和高阶 Moyal 导数都记录着经典概率论无法容纳的量子结构。TWA 所做的是在明确条件下保留初态的相空间涨落和领先的半经典动力学，同时舍弃一部分后续量子干涉。因此，任何应用章节都同时讲授计算方法和失效诊断。

全书由简单到复杂：单模谐振子建立符号和直觉，Kerr 振子暴露截断的代价，Bose--Hubbard 二聚体连接多模量子动力学，离散玻色场引出实际数值问题，最后以双分量凝聚体淬火说明如何构造 Bogoliubov 初态、传播随机场并诊断密度序与相干序。涉及数值的图表均由配套程序实际生成，而不是根据预期结论绘制。

这是一本方法教材，而不是对某个具体体系预设结论的论证。若模型、参数或有限尺寸证据不足，本书将结论停在证据能够支持的位置。
